\vspace{-0.5em}
\section{Related Work}\label{sec:rw}

Quantifying trust and uncertainty in neural networks is central in safety-critical domains, where incorrect yet confident predictions can have severe consequences; under data poisoning, a model may score well on standard metrics while having learned erroneous patterns. Diverse frameworks model and propagate uncertainty and trust, yet significant limitations remain, as we point out next.
\vspace{-0.5em}
\subsection{Uncertainty Quantification in Neural Networks}

Uncertainty quantification estimates predictive reliability by modeling epistemic and aleatoric uncertainty~\cite{kendall2017uncertainties}, through Bayesian neural networks~\cite{neal1996bayesian,blundell2015weightuncertaintyneuralnetworks,pmlr-v80-depeweg18a}, Monte Carlo dropout~\cite{gal2016dropoutbayesianapproximationrepresenting}, and ensembles~\cite{lakshminarayanan2017simplescalablepredictiveuncertainty}. Yet calibration remains difficult in complex models~\cite{Gawlikowski2023}, confidence degrades under dataset shift~\cite{begoli2019need,2506.07461,1906.02530}, and post-hoc calibration such as temperature scaling~\cite{guo2017calibrationmodernneuralnetworks} acts only at the output layer and assumes clean data, so it can still mislead when inputs or training data are corrupted.
\vspace{-0.5em}
\subsection{Subjective Logic Approaches to Trust}

Evidential deep learning (EDL)~\cite{sensoy2018evidentialdeeplearningquantify} applies SL principles by representing class predictions as subjective opinions parameterized through a Dirichlet distribution, jointly predicting outcomes and quantifying confidence so as to distinguish low-confidence predictions from high-uncertainty regions such as out-of-distribution inputs. Although effective, EDL assumes clean data and provides no input-level trust assessment. A calibration-based approach~\cite{11124121} instead clusters model outputs into subjective opinions to derive per-prediction trust scores without accessing internal parameters; while simple and model-agnostic, it likewise assumes a trustworthy dataset and neglects input evidence.

Two further lines relate directly to our evaluation. Input- and feature-space detectors flag distribution violations without touching training dynamics: Mahalanobis distances on hidden features~\cite{lee2018mahalanobis}, deep nearest-neighbor distances~\cite{sun2022knnood}, energy scores~\cite{liu2020energy}, and conformal prediction~\cite{angelopoulos2023conformal} wrap point predictions with distribution-level evidence; \cref{exp:4} includes the first three as baselines and quantifies what trust propagation adds beyond them. Data-centric methods attribute unreliability to training points: influence functions~\cite{koh2017influence} and confident learning~\cite{northcutt2021confidentlearning} estimate which labels are erroneous and grade corruption severity, and we compare against such signals in the training-quality audit. PaTAS differs from both lines by composing input-, parameter-, and path-level trust into a single propagating calculus with an interpretable opinion output, rather than producing a task-specific score.

The DeepTrust framework by Cheng et~al.~\cite{10.3389/frai.2020.00054} adopts a white-box perspective, integrating dataset evidence during training to assess global model trustworthiness. Although holistic, its use of SL operators lacks algebraic consistency: it maps neural-network addition to SL fusion and multiplication to SL opinion multiplication, although fusion operates over agents' opinions whereas multiplication operates over variables. Moreover, its trust backpropagation is specified only at a single-layer level, leaving the extension to deeper architectures unclear despite empirical evaluations on complex networks. PaTAS addresses these limitations through a coherent propagation mechanism defined layer-wise, hence at arbitrary depth, with consistency properties established analytically, convergence proven for the idealized fixed-revision iteration, and training-time stability assessed empirically (\cref{sec:ptastheory}).
\vspace{-0.5em}
\subsection{Trust and Uncertainty Propagation}

Work on uncertainty propagation improves accuracy and efficiency: sampling-free variance propagation for dropout networks~\cite{MAE2021394}, Gaussian-mixture propagation of input uncertainty~\cite{pmlr-v202-monchot23a}, and extensions to multi-layer perceptrons~\cite{astudillo2011propagation}. Beyond neural models, the Appleseed model propagates trust in social networks via spreading activation~\cite{ziegler2004spreading}, showing the value of treating trust as a structural, context-dependent quantity.

These approaches show growing interest in uncertainty and trust modeling, yet they typically treat reliability as a property of the output layer or of the model as a whole, and none propagates trust jointly from input, training data, and parameters. PaTAS differs in three specific respects rather than in its use of Subjective Logic alone. First, it defines a trust computation that \emph{structurally mirrors} the network: each neuron has a Trust Node, so trust is propagated along the same topology as the prediction rather than estimated at the output. Second, it derives \emph{parameter-level} trust during training from gradient evidence conditioned on label trust, which existing evidential and calibration methods do not model. Third, its Inference-Path Trust Assessment yields \emph{per-instance, path-specific} trust at deployment. The individual ingredients, Subjective Logic and evidential reasoning are known; the contribution is their composition into a single mechanism that assigns and propagates trust at input, parameter, and path granularity simultaneously, together with the consistency properties this mechanism is proven to satisfy.